\chapter{Asset Bridging}\label{app:bridging}

This appendix specifies the concrete data structures and abstract algorithms of external asset bridging introduced in Sec.~\ref{sec:bridged-assets}. The presentation is chain-independent: the external settlement chain is abstracted as a contract with persistent storage and a succinct-proof verifier (an account-based EVM chain is the motivating instance). Cryptographic primitives are treated as black boxes: the hash $H$ is $\mathsf{SHA\text{-}256}$ over deterministic CBOR (Appendix~\ref{app:encodings}), the authenticated set reuses the sparse Merkle tree of Appendix~\ref{app:rsmt}, and the proof system is specified by its interface only (Sec.~\ref{app:bridging-primitives}).

All multi-field hash inputs are CBOR arrays in the stated field order; all structures use deterministic CBOR and MUST be consumed completely, rejecting non-canonical encodings, wrong tags, wrong field counts or types, and trailing bytes, as elsewhere in this specification.

\section{Cryptographic Primitives}\label{app:bridging-primitives}

\paragraph{Succinct proof system.}
A proof system for an NP relation $\mathfrak{R}$ is a pair of algorithms over a public statement $\mathsf{x}$ and a private witness $\mathsf{w}$:
\[
\Pi \gets \mathsf{Prove}(\mathsf{x},\mathsf{w}), \qquad
\mathsf{Verify}(\mathsf{vk},\mathsf{x},\Pi) \in \bool,
\]
with a verification key $\mathsf{vk}$ that pins the relation. It MUST satisfy:
\begin{itemize}
  \item \textbf{Completeness}: if $(\mathsf{x},\mathsf{w}) \in \mathfrak{R}$ then $\mathsf{Verify}(\mathsf{vk},\mathsf{x},\mathsf{Prove}(\mathsf{x},\mathsf{w}))=1$.
  \item \textbf{Knowledge soundness}: a prover that makes $\mathsf{Verify}$ accept knows a witness $\mathsf{w}$ with $(\mathsf{x},\mathsf{w}) \in \mathfrak{R}$, except with negligible probability.
  \item \textbf{Succinctness}: the cost of $\mathsf{Verify}$ is constant in $|\mathsf{w}|$ (in particular independent of token-history length and batch size).
  \item \textbf{Recursive composability}: a proof of $\mathfrak{R}$ can be verified \emph{inside} a relation, at a cost independent of the statement it establishes. Required by token compression (Sec.~\ref{sec:compression-relation}, check~5) and by the attestation shortcut of the return relation below; not required by a deployment that expands tokens inline.
\end{itemize}
The external vault hard-codes $\mathsf{vk}$; $\mathsf{x}$ is supplied on-chain and $\mathsf{w}$ never leaves the prover.

\paragraph{Instantiation.}
More than one family satisfies the interface above, and the choice is a deployment parameter rather than a property of any relation. The default profile realizes both relations as programs for a general-purpose zero-knowledge virtual machine, which is what allows irregular control flow, canonical CBOR decoding and sorting to be written as ordinary code and to share source with the verifiers outside the relation. Its native output is a \emph{recursive STARK}, used throughout the Unicity ecosystem: directly recursion-friendly, no trusted setup, constant verification cost, at the price of a proof length on the order of megabytes. That is the form $\pi_\mathsf{tok}$ takes (Sec.~\ref{sec:compression-relation}), the form the compression relation verifies in-circuit, and the form every relying party inside Unicity verifies.

A pairing-based wrap --- Groth16 over a pairing-friendly curve --- is applied at exactly one boundary: the return path specified in this appendix, where the verifier is a contract paying per calldata byte and per operation. The wrap is a final step of the return prover over a statement the recursive proof has already established, and it is produced on demand rather than carried. No proof circulating inside Unicity is therefore wrapped, which leaves one proof form in the token format, one $\mathsf{vk}_\mathsf{tok}$, and no attestation that is verifiable but not foldable. Where the return relation discharges an attestation by verifying its $\pi_\mathsf{tok}$ in-circuit (Sec.~\ref{app:bridging-relation}), it consumes that unwrapped form.

A deployment MAY substitute another wrapping scheme, or another proof system entirely. The requirements are that the vault's pinned $\mathsf{vk}$ and the wrapping scheme are chosen together, that a change of either is a vault redeployment (Appendix~\ref{app:bridging-upgrades}), and that whatever is used inside Unicity satisfies recursive composability.

\paragraph{Nullifier accumulator.}
An authenticated set $\mathcal{D}$ over $\bytes{32}$ is summarized by a root $A = \mathsf{Root}(\mathcal{D}) \in \bytes{32}$, with empty-set root $A_\emptyset$, and operations
\[
\begin{aligned}
&\mathsf{NonMember}(A,\eta) \to w, &&\mathsf{VerifyNonMember}(A,\eta,w) \in \bool,\\
&\mathsf{Insert}(A,\eta,w) \to A',
\end{aligned}
\]
where $w$ is a membership-state witness. Required properties: $\mathsf{VerifyNonMember}(A,\eta,w)=1$ certifies $\eta \notin \mathcal{D}$ under root $A$; $\mathsf{Insert}$ returns the root of $\mathcal{D}\cup\{\eta\}$; and it is infeasible to produce a valid non-membership witness for an element already inserted. The set is realizable by the sparse Merkle tree of Appendix~\ref{app:rsmt} keyed by $\eta$ (its inclusion and non-inclusion certificates implement $\mathsf{VerifyMember}$/$\mathsf{VerifyNonMember}$).

\paragraph{Batch insertion is ordered.}
Inserting a sequence $\langle\eta_1,\ldots,\eta_B\rangle$ applies $\mathsf{Insert}$ one element at a time, threading the root. The witness $w_k$ for $\eta_k$ MUST be valid against the intermediate root \emph{after} $\eta_1,\ldots,\eta_{k-1}$ have been inserted, not against the starting root. A witness builder therefore simulates the batch insertions in order; a witness computed against $A_\mathsf{old}$ is valid only for $\eta_1$.

\section{Configuration and Domain Separators}\label{app:bridging-config}

A bridge deployment fixes a configuration
\[
\mathcal{C} = (\chi,\; a_\mathsf{V},\; a_\mathsf{A},\; \mathsf{ty},\; \mathsf{aid},\; d_\mathsf{lock},\; d_\mathsf{nul},\; t_\mathsf{reason}),
\]
where $\chi$ is the external chain identifier, $a_\mathsf{V}$ the vault address, $a_\mathsf{A}$ the external asset (token contract) address, $\mathsf{ty}$ the bridged token type, $\mathsf{aid}$ the asset identifier read by $\Call{Assets}{\mathsf{ty},\auxd'}$ (Sec.~\ref{sec:asset-structure}), $d_\mathsf{lock}$ and $d_\mathsf{nul}$ are the lock-digest and nullifier domain separators, and $t_\mathsf{reason}$ is the CBOR semantic tag of the return reason. The \emph{configuration hash} is
\[
h_\mathsf{cfg} = H_\mathsf{cfg}(\mathcal{C})
= H(\mathsf{CFG\_DOMAIN},\; \chi,\; a_\mathsf{V},\; a_\mathsf{A},\; \mathsf{ty},\; \mathsf{aid},\; d_\mathsf{lock},\; d_\mathsf{nul},\; t_\mathsf{reason}).
\]
The vault MUST derive $h_\mathsf{cfg}$ and its custody asset $a_\mathsf{A}$ from the \emph{same} configuration value at deployment, so that the asset the proof validates against and the asset the vault pays out cannot diverge. $\mathsf{CFG\_DOMAIN}$ is fixed by this specification; $d_\mathsf{lock}$, $d_\mathsf{nul}$, and $t_\mathsf{reason}$ are deployment-pinned fields. In algorithms below, $\mathsf{LOCK\_DOMAIN}$, $\mathsf{NUL\_DOMAIN}$, and $\mathsf{REASON\_TAG}$ denote the corresponding fields of the already-bound $\mathcal{C}$.

Note that $\mathcal{C}$ pins \emph{which} asset identifier is bridged but not how $\Call{Assets}{\mathsf{ty},\auxd'}$ decodes it (Sec.~\ref{sec:asset-structure}) --- that decoding is fixed by the token type $\mathsf{ty}$ itself and must be identical across every independent implementation that mints, backs, or returns the bridged token.

\section{Bridge-In: Lock, Record, and Backing Proof}\label{app:bridging-lock}

Bridge-in escrows the external asset and mints a Unicity token whose genesis backing reason ties it to that escrow. This section gives the lock event, the digest the vault retains for return-time backing, the backing-reason structure, the external-finality black box, and the verification algorithm a recipient runs.

\subsection{Lock Operation, Event, and Record}\label{app:bridging-lock-op}

The vault's lock operation escrows the deposited asset into pooled custody, assigns a monotone \emph{lock nonce} $n$, and emits a \emph{lock event}
\[
L_\mathsf{lock} = (n,\; a_\mathsf{src},\; v,\; \mathsf{id},\; h_\mathsf{rcpt}),
\]
where $a_\mathsf{src}$ is the depositor, $v$ the locked amount, $\mathsf{id}$ the identifier of the token to be funded, and $h_\mathsf{rcpt} = H(\mathsf{CBOR}(\predi'_0))$ the commitment to that token's first-owner predicate. Simultaneously the vault records a commitment for return-time use. Define the \emph{lock record}
\[
\mathcal{K} = (a_\mathsf{A},\; \mathsf{ty},\; \mathsf{aid},\; v,\; \mathsf{id},\; h_\mathsf{rcpt})
\]
and store only its \emph{lock digest}
\[
d \;=\; H(\mathsf{LOCK\_DOMAIN},\; \chi,\; a_\mathsf{V},\; n,\; \mathcal{K}),
\qquad \mathsf{lockDigest}[n] \gets d .
\]
The fields $a_\mathsf{A},\mathsf{ty},\mathsf{aid}$ are deployment constants from $\mathcal{C}$; $v,\mathsf{id},h_\mathsf{rcpt}$ come from the lock. Binding $\mathsf{id}$ into both the lock and the genesis ties one lock to exactly one token (token-identifier uniqueness, Sec.~\ref{sec:mint-tx}), and binding $h_\mathsf{rcpt}$ ensures only the designated recipient owns the minted token. The digest is the only return-relevant state the vault retains; it lets the return relation (Sec.~\ref{app:bridging-relation}) re-establish backing without any external query.

\subsection{Backing Reason}\label{app:bridging-backing-reason}

The genesis mint reason (Sec.~\ref{sec:mint-justification}) of a bridged token is a CBOR array under tag $\mathsf{BACKING\_TAG}$ that locates the lock and pins the trust anchors:

\begin{longtable}{clp{8.4cm}}
\textbf{Index} & \textbf{Field} & \textbf{Description} \\
\hline
\endhead
0 & $\mathsf{ver}$ & version, $=1$ \\
1 & $\chi$ & external chain identifier \\
2 & $a_\mathsf{V}$ & vault address \\
3 & $a_\mathsf{A}$ & asset address \\
4 & $n$ & lock nonce (locator of the lock event) \\
\end{longtable}

The reason body carries only a reference to the lock and the trust anchors; the security-relevant values --- token identifier, recipient, and amount --- are read from the certified genesis and the lock event and checked against each other, never trusted from the reason body. A deployment MAY additionally carry an external-chain inclusion proof of the lock in the reason (Sec.~\ref{app:bridging-finality}). The structure of the backing reason and its verifier are part of the token type $\mathsf{ty}$'s definition.

\subsection{External Finality}\label{app:bridging-finality}

Confirming the lock is the one step Unicity cannot internalize, since it concerns another chain's state. It is abstracted as a black-box oracle
\[
\Call{LockFinal}{\chi, a_\mathsf{V}, n} \;\to\; (\mathsf{ok},\; a_\mathsf{A},\; v,\; \mathsf{id},\; h_\mathsf{rcpt}),
\]
returning $\mathsf{ok}=1$ together with the committed fields of $L_\mathsf{lock}$ iff a lock with nonce $n$ was emitted by vault $a_\mathsf{V}$ on chain $\chi$ and is final there. Two realizations satisfy the contract:
\begin{enumerate}
  \item \textbf{Live query}: the verifier queries a $\chi$ node and requires at least $K$ confirmations, with $K$ chosen for $\chi$'s reorg resistance.
  \item \textbf{Carried inclusion proof}: the reason carries a $\chi$ block header and a log/receipt inclusion proof, verified against a pinned $\chi$ trust anchor; finality follows from the header depth or the anchor's finalization rule.
\end{enumerate}
Both lie outside Unicity's trust base $\mathcal{T}$ and are the bridge's only external-chain assumption. The choice trades liveness and connectivity (live query) against self-containment (carried proof) and is a per-deployment parameter.

\subsection{Backing Verification Algorithm}\label{app:bridging-backing}

Algorithm~\ref{alg:bridge-backing} is the token-type mint-reason check (Sec.~\ref{sec:mint-tx}, step~5) for a bridged genesis; it runs after the generic mint-transaction verification, against the pinned configuration $\mathcal{C}$ (Sec.~\ref{app:bridging-config}).

\begin{algorithm}[!htbp]
\caption{Backing Mint-Reason Verification}\label{alg:bridge-backing}
\begin{algorithmic}[0]
\Function{VerifyBackingReason}{$D_\mathsf{mint},\mathcal{T},\mathcal{C}$}
  \State $(\mathsf{ver},\chi',a'_\mathsf{V},a'_\mathsf{A},n) \gets \Call{DecodeBackingReason}{D_\mathsf{mint}.\mathsf{reason}}$
  \State ensure($\mathsf{ver} = 1$)
  \State ensure($(\chi',a'_\mathsf{V},a'_\mathsf{A}) = (\mathcal{C}.\chi,\mathcal{C}.a_\mathsf{V},\mathcal{C}.a_\mathsf{A})$) \Comment{trust anchors}
  \State ensure($D_\mathsf{mint}.\mathsf{ty} = \mathcal{C}.\mathsf{ty}$ \textbf{and} $D_\mathsf{mint}.\alpha = \mathcal{T}.\alpha$)
  \State $(\mathsf{ok},a''_\mathsf{A},v,\mathsf{id}_L,h_\mathsf{rcpt}) \gets \Call{LockFinal}{\chi',a'_\mathsf{V},n}$ \Comment{external finality oracle}
  \State ensure($\mathsf{ok} = 1$ \textbf{and} $a''_\mathsf{A} = a'_\mathsf{A}$)
  \State $\mathsf{id} \gets H(D_\mathsf{mint}.\mathsf{salt},D_\mathsf{mint}.\alpha)$
  \State ensure($\mathsf{id}_L = \mathsf{id}$) \Comment{lock funds this exact token}
  \State ensure($h_\mathsf{rcpt} = H(\mathsf{CBOR}(D_\mathsf{mint}.\predi'))$) \Comment{lock binds this recipient}
  \State $v_g \gets$ value of $D_\mathsf{mint}.\auxd'$ for $\mathcal{C}.\mathsf{aid}$ via $\Call{Assets}{D_\mathsf{mint}.\mathsf{ty},D_\mathsf{mint}.\auxd'}$
  \State ensure($v = v_g$) \Comment{minted value equals locked amount}
  \State \Return $1$
\EndFunction
\end{algorithmic}
\end{algorithm}

\subsection{Bridge-In Security}\label{app:bridging-in-security}

\begin{description}
\item[Backing at mint time.] An accepted bridged genesis implies a final lock on $\chi$ committing to the same identifier, recipient, and amount (Alg.~\ref{alg:bridge-backing}); every bridged token therefore first circulates fully backed.
\item[Uniqueness.] Exactly one genesis exists per $\mathsf{id}$ (Sec.~\ref{sec:request-validation}), and the lock fixes $\mathsf{id}$; hence each lock funds one token, and no second genesis can claim the same lock.
\item[Backing at return time.] The return relation does not re-invoke $\Call{LockFinal}{}$. It reconstructs $\mathcal{K}$ from the certified genesis and $\mathcal{C}$, recomputes $d$, and the vault confirms $d = \mathsf{lockDigest}[n]$ (Sec.~\ref{app:bridging-settlement}). By collision resistance this matches only the lock the vault actually recorded, so custody is released solely against a genuine deposit, with no external query at return time.
\end{description}

\section{Return Reason}\label{app:bridging-reason}

The return reason is a CBOR array under tag $\mathsf{REASON\_TAG}$:

\begin{longtable}{clp{8.4cm}}
\textbf{Index} & \textbf{Field} & \textbf{Description} \\
\hline
\endhead
0  & $\mathsf{ver}$        & version, $=1$ \\
1  & $\chi$                & external chain identifier \\
2  & $a_\mathsf{V}$        & vault address \\
3  & $a_\mathsf{A}$        & asset address \\
4  & $\mathsf{ty}$         & bridged token type \\
5  & $\mathsf{aid}$        & asset identifier \\
6  & $a_\mathsf{r}$        & external recipient of the release \\
7  & $v$                   & gross amount; equals the burned token's value for $\mathsf{aid}$ \\
8  & $a_\mathsf{f}$        & fee recipient ($\bot$ disables the fee) \\
9  & $v_\mathsf{f}$        & fee amount, $v_\mathsf{f} \le v$ \\
10 & $t_d$                 & deadline; gates the fee only (Sec.~\ref{app:bridging-settlement}) \\
\end{longtable}

Let $b_\mathcal{R}$ be the deterministic CBOR encoding of $\mathcal{R}$ under $\mathsf{REASON\_TAG}$. The holder burns the source token (Sec.~\ref{sec:burn-transaction}) with terminal auxiliary data $b_\mathcal{R}$ and recipient predicate $\predi_\mathsf{burn}(H(b_\mathcal{R}))$. The return reason is therefore part of the certified burned token, not a separate witness supplied by the prover. Field $0$ pins the reason version, fields $1$--$5$ MUST equal $(\mathcal{C}.\chi,\mathcal{C}.a_\mathsf{V},\mathcal{C}.a_\mathsf{A},\mathcal{C}.\mathsf{ty},\mathcal{C}.\mathsf{aid})$, and fields $6$--$10$ are the per-return release parameters. If $a_\mathsf{f}=\bot$, the fee is inactive regardless of $v_\mathsf{f}$; settlement pays the full value to $a_\mathsf{r}$ and performs no transfer to $\bot$. Every field is public by nature, since the external release is a public event; the private part of a return is the token's transfer history, which never leaves the witness.

\section{Nullifier}\label{app:bridging-nullifier}

Let $\mathsf{BURN\_DOMAIN}$ be the fixed byte string \texttt{unicity-burn-transition:v1}. Let the burn be the final certified transfer of the source token. It spends the state $(\predi_{n-1},\sthash_{n-1})$ and its transaction data hash is $\txhash_\mathsf{b}=H(T_\mathsf{b}.D)$. The certified Unicity leaf for that burn is keyed by
\[
\mathsf{sid}_\mathsf{b} = H(\predi_{n-1},\sthash_{n-1}).
\]
The burn transition identifier and \emph{nullifier} are
\[
\mathsf{btid}_\mathsf{b} = H(\mathsf{BURN\_DOMAIN},\; \mathsf{sid}_\mathsf{b},\; \txhash_\mathsf{b}),
\qquad
\eta \;=\; H(\mathsf{NUL\_DOMAIN},\; h_\mathsf{cfg},\; \mathsf{btid}_\mathsf{b}).
\]
The relation MUST derive $\mathsf{sid}_\mathsf{b}$ and $\txhash_\mathsf{b}$ from the final certified burn itself; they are not independent witness inputs. $\mathsf{sid}_\mathsf{b}$ is recorded at most once by the Unicity Service (Sec.~\ref{sec:request-validation}), so $\eta$ is globally unique; the certified $\txhash_\mathsf{b}$ binds it to the exact return reason, and $h_\mathsf{cfg}$ separates deployments so the same burn never collides across independent bridges.

\section{Shared Anchor for Inclusion Checking}\label{app:bridging-anchor}

Many certified transactions can be authenticated against a single verified signature. This section specifies the primitive as used by the bridge-return relation below and by token compression (Sec.~\ref{sec:token-compression}).

\subsection{The Anchor}\label{app:anchor-level}

An \emph{anchor} is a Unicity Seal $C^{\mathsf{r}*}$, and the value it contributes is a single scalar plus a size: the \emph{anchor accumulator} $a^* = C^{\mathsf{r}*}.a_r$ and its size $N^* = C^{\mathsf{r}*}.n_r + 1$ (Sec.~\ref{sec:root-history}), which together commit every Unicity Tree root of rounds $0,\ldots,C^{\mathsf{r}*}.n_r$. $N^*$ is derived from the signed round number, so it is authenticated by the same signature; it must be carried explicitly because the accumulator's shape is a function of its size, and a verifier that does not fix the size does not fix the structure it is checking against. The seal's own root $C^{\mathsf{r}*}.r$ is not used separately; it is the last element $a^*$ commits, and treating it as a second anchor component only creates a value that has to be authenticated on its own.

The natural-looking alternative --- anchoring at a shard's SMT root $UC^*.\mathcal{IR}.h$ --- does not work once a partition has more than one shard, and the reason is structural rather than a matter of tuning. A state identifier $\mathsf{sid} = H(\predi, \sthash)$ is a hash, so the successive states of one token are distributed across shards pseudo-randomly by $f_\mathcal{SH}$ (Sec.~\ref{sec:sharding}). A shard root commits only its own key range, so a shard-level anchor can never cover a leaf in another shard: a token with $m$ transitions over $S$ shards would need $\min(m, S)$ anchor certificates, each carrying a full BFT quorum signature, and no amount of re-anchoring would reduce that number. Nor can it be repaired by asking the shards to certify in step --- shards run asynchronously and no mechanism synchronizes their rounds.

Anchoring one level up dissolves both problems. The seal signs $r$, the root of the Unicity Tree, which commits every shard of every partition through the shard trees (Sec.~\ref{VerifyUnicityCert}), and the shard tree carries forward each shard's last certified input record when that shard has no input to certify in a round (Sec.~\ref{se:shard-tree}). A seal at BFT round $n$ therefore commits \emph{some} certified SMT root for \emph{every} shard, whatever the shards' relative progress. Shard asynchrony becomes irrelevant: one anchor covers leaves from any set of shards, and the only requirement is that the anchor be recent enough that each shard's committed input record post-dates the leaves taken from it.

\subsection{Anchored Inclusion}

Anchoring is two-staged: one batched step authenticates the Unicity Tree roots of every round a proof draws leaves from, and a per-leaf step then works against one of those authenticated roots. The split is not presentational --- it is what makes the root history cost proportional to the number of distinct \emph{rounds} rather than to the number of leaves, and a verifier that checks a separate history path per leaf gets none of the saving claimed below.

\paragraph{Round table.} A proof declares the distinct rounds it draws leaves from,
\[
\mathcal{R} = \langle (n_1, \rho_1), \ldots, (n_d, \rho_d)\rangle, \qquad n_1 < \cdots < n_d ,
\]
where $\rho_j$ is the Unicity Tree root of round $n_j$. The whole table is authenticated in one step:
$\Call{VerifyRounds}{a^*, N^*, \mathcal{R}, \Psi} = 1$ iff the indices are strictly increasing, $n_d < N^*$, and $\Psi$ verifies $\mathcal{R}$ against $(a^*, N^*)$ as a batched root history proof (Sec.~\ref{sec:root-history}). Verification work is $O(d\log(1 + N^*/d) + \log N^*)$: the auxiliary path material shrinks as $d$ grows and the paths converge, but the $d$ table entries must each be consumed, so the $\Omega(d)$ term remains. At $d = N^*$ this is $\Theta(N^*)$, as it must be.

\paragraph{Leaf.} Against the authenticated table, a leaf $L = (\mathsf{sid}, v)$ is established by
\[
\pi^\mathsf{an} = (j,\, C^\mathsf{inc},\, \mathcal{IR},\, h_t,\, C^\mathsf{shard},\, C^\mathsf{uni},\, \mathcal{CD}_\beta),
\]
accepted by $\Call{VerifyAnchoredLeaf}{\alpha, \mathcal{R}, L, \pi^\mathsf{an}}$ iff, writing $\rho \gets \mathcal{R}[j].\rho$:
\begin{enumerate}
  \item $\Call{rsmt\_verify\_inclusion}{C^\mathsf{inc}, \mathcal{IR}.h, \mathsf{sid}, v} = 1$ \quad (the leaf is in the shard's SMT);
  \item $\mathcal{CD}_\beta.\alpha = \alpha$ and $H(\mathcal{CD}_\beta) = C^\mathsf{uni}.\mathsf{dhash}$ \quad (\emph{network binding} and authentic partition description);
  \item $\sigma \gets C^\mathsf{shard}.\sigma$;\; $\sigma \in \mathcal{CD}_\beta.\mathcal{SH}$ and $f_{\mathcal{CD}_\beta.\mathcal{SH}}(\mathsf{sid}) = \sigma$ \quad (\emph{right-shard binding});
  \item $\mathsf{CompUnicityTreeCert}(C^\mathsf{uni}, \mathsf{CompShardTreeCert}(C^\mathsf{shard}, \mathcal{IR}, h_t)) = \rho$ \quad (the leaf's round is the one claimed).
\end{enumerate}

Checks 2 and 3 are the two bindings that ordinary inclusion verification performs and that the certificate chain does not imply on its own. $\mathsf{VerifyUnicityCert}$ deliberately does \emph{not} bind a state identifier to the shard named in $C^\mathsf{shard}$ (Sec.~\ref{VerifyUnicityCert}); that binding is supplied by $\mathsf{VerifyInclusionProof}$ (Sec.~\ref{VerifyInclusionProof}), which also requires $\mathcal{CD}_\beta.\alpha = \mathcal{T}.\alpha$. Anchored verification must reproduce both, or it is not equivalent to the ordinary path: without check 3 it would accept a leaf recorded by a shard not responsible for its key --- the exact equivocation the per-shard key-space invariant exists to prevent (Sec.~\ref{sec:partitions-shards}) --- and without check 2 it would accept a partition description belonging to another network.

Every round root is established through the accumulator, with no shortcut for the anchor's own round. Admitting such a branch would mean the anchor consists of two values, of which one --- the accumulator --- would then have no reason to be authenticated wherever the pair is not read straight off a verified seal. Since an accumulator authorizes arbitrary historical roots, an unauthenticated one authorizes arbitrary fabricated roots. One value, one batched table, no branch.

With it, one recent seal authenticates the roots of all earlier rounds, so a history of any length reduces to one signature verification, and the relying party never needs the trust base of an earlier epoch. Authority flows from the one seal it verifies under its \emph{current} trust base, and an old seal from an abandoned fork cannot be substituted, because its root is not in the accumulator of the chain the relying party is on. Token compression (Sec.~\ref{sec:token-compression}) extends the same property to provenance depth, subject to how each fold reconciles the anchor of the attestation it absorbs.

Cost per leaf is one SMT path plus $|\sigma| + O(\log |\mathcal{C}|)$ hashes for the shard and unicity trees; cost per proof is the round table and one seal verification, the latter being the only BFT signature check involved.

This is not free asymptotically, and the trade should be stated plainly. Anchoring at a single shard root cost $O(1)$ certification evidence per leaf but, as shown above, cannot span shards at all. Anchoring at a seal costs $O(d\log(1 + N^*/d) + \log N^*)$ for $d$ distinct rounds instead. What it replaces is one \emph{Unicity Certificate verification} per leaf --- a BFT quorum signature check --- with hash compressions shared across a round, against a single signature check for the whole proof. For any realistic $N^*$ the added hashing is far cheaper than the signature verifications removed, but the added factor is real and does not vanish.

\subsection{Reference Time from an Anchor}\label{app:anchor-time}

The reference time of a transition is not read from the round the leaf is anchored at. It is committed by the leaf value itself, $v = H(\txhash, \tau)$, fixed by the BFT Core when the round was certified (Sec.~\ref{sec:us-processing}, Sec.~\ref{sec:derived-leaf-values}). A prover witnesses $\tau$, check~1 above binds it by reconstructing $v$, and a time-dependent predicate --- $\predi_\mathsf{tlock}$, $\predi_\mathsf{htlc}$ --- is then evaluated in-circuit exactly as outside it.

Anchoring is therefore uniform over predicate types. Any leaf may be established against any round at or after its own, so every leaf of a derivation can be re-presented against a recent round (Sec.~\ref{sec:leaf-representation}) and no class of transition obliges a prover to obtain a root history path for a distant round. Were $\tau$ instead read off the anchored round's input record, a time-dependent transition would have to be presented at its own round and would carry a retention dependency the others do not (Sec.~\ref{sec:root-history-pruning}); committing $\tau$ to the leaf removes that asymmetry rather than managing it.

\section{Public Statement and Public Calldata}\label{app:bridging-public}

The public statement committed by the proof and checked by the vault is
\[
\mathsf{x} = (h_\mathsf{cfg},\; h_\mathcal{T},\; A_\mathsf{old},\; A_\mathsf{new},\; \mathsf{rr},\; \mathsf{lr},\; B,\; V),
\]
with $h_\mathcal{T} = H(\mathcal{T})$, $A_\mathsf{old}/A_\mathsf{new}$ the accumulator root transition, $B$ the batch size, $V$ the total gross amount, and $\mathsf{rr}, \mathsf{lr}$ commitments to the calldata below. The vector commitment $\Call{Commit}{}$ is fixed by the verification key and the vault implementation; it is not chosen per batch.

A \emph{release leaf} carries everything the vault must execute and emit:
\[
\ell = (\eta,\; a_\mathsf{r},\; v,\; a_\mathsf{f},\; v_\mathsf{f},\; t_d),
\qquad \mathsf{rr} = \Call{Commit}{\langle \ell_1,\ldots,\ell_B \rangle}.
\]
The nullifier $\eta$ is included in the leaf, and emitted on settlement, so any party can reconstruct the off-chain accumulator and replace a stalled prover. A \emph{lock reference} is $(n,d)$, and $\mathsf{lr} = \Call{Commit}{\cdot}$ over the distinct lock references sorted by $n$.

\section{Batch Proving Relation}\label{app:bridging-relation}

Algorithm~\ref{alg:bridge-relation} defines $\mathfrak{R}$: the relation holds iff every check passes. The witness $\mathsf{w}$ comprises the bridge configuration $\mathcal{C}$, the trust base $\mathcal{T}$, the anchor seal $C^{\mathsf{r}*}$, the round table $\mathcal{R}$ with its batched proof $\Psi$, and for each burn $i$: the source token $\mathcal{L}_i$, its anchored inclusion proofs, the lock record(s) backing its genesis, and the accumulator non-membership witness $w_i$. The first check binds the private configuration to the public statement by recomputing $h_\mathsf{cfg}$ exactly as in Sec.~\ref{app:bridging-config}; $\mathcal{C}$ is then safe to use for all per-burn checks. $\Call{VerifyToken}{}$ (Sec.~\ref{sec:token-verification}) is invoked in \emph{anchored} mode: wherever it would verify a certified transaction's own Unicity Certificate, it instead verifies an anchored leaf against the round table $\mathcal{R}$ (Sec.~\ref{app:bridging-anchor}); its reference-form split-reason checks (Sec.~\ref{sec:split-verification}) are unchanged and supply value conservation, with each source resolved from the witness rather than from an evidence set. Reference time is available from the anchored path (Sec.~\ref{app:anchor-time}), so time-dependent predicates are admitted. The genesis backing verifier is \emph{structural}: it checks the binding to the lock record and exports the lock reference as a public obligation, leaving the existence of the lock to the vault's recorded digest (Sec.~\ref{app:bridging-in-security}). The external finality of the lock was already established at mint time by the bridge-in verifier (Sec.~\ref{app:bridging-backing}), so it need not be re-checked inside the proof.

A deployment that already verifies token attestations (Sec.~\ref{sec:token-compression}) may discharge $\Call{VerifyToken}{}$ here by verifying an attestation's $\pi_\mathsf{tok}$ in-circuit instead of expanding the token, reconciling its anchor against $a^*$ by either of the routes $\mathfrak{R}_\mathsf{tok}$ admits (Sec.~\ref{sec:compression-relation}, check~5); the two relations share the same anchoring primitive by construction. Taking the seal route requires the trust base to be reachable inside this relation, which it is, since $\mathcal{T}$ is already a witness here and $h_\mathcal{T}$ is already part of the public statement.

\begin{algorithm}[!htbp]
\caption{Batch Return Relation $\mathfrak{R}(\mathsf{x},\mathsf{w})$}\label{alg:bridge-relation}
\begin{algorithmic}[0]
\Function{BridgeReturnRelation}{$\mathsf{x},\mathsf{w}$}
  \State $\mathcal{C} \gets \mathsf{w}.\mathcal{C}$
  \State ensure($H_\mathsf{cfg}(\mathcal{C}) = \mathsf{x}.h_\mathsf{cfg}$)
  \State $(\chi,a_\mathsf{V},a_\mathsf{A},\mathsf{ty},\mathsf{aid},d_\mathsf{lock},d_\mathsf{nul},t_\mathsf{reason}) \gets \mathcal{C}$
  \State ensure($H(\mathcal{T}) = \mathsf{x}.h_\mathcal{T}$)
  \State ensure($\mathsf{VerifyUnicitySeal}(C^{\mathsf{r}*}.r, C^{\mathsf{r}*},\mathcal{T}) = 1$);\quad $(a^*, N^*) \gets (C^{\mathsf{r}*}.a_r,\; C^{\mathsf{r}*}.n_r + 1)$
  \State ensure($\Call{VerifyRounds}{a^*, N^*, \mathcal{R}, \Psi} = 1$) \Comment{one batched round table for the whole batch}
  \State $A \gets \mathsf{x}.A_\mathsf{old}$;\quad $V \gets 0$;\quad $\mathsf{leaves} \gets \langle\rangle$;\quad $\mathsf{lref} \gets \langle\rangle$
  \For{$i \gets 1 \ldots \mathsf{x}.B$}
    \State \textbf{Finality:} ensure($\Call{VerifyToken}{\mathcal{L}_i,\mathcal{T}} = 1$ in anchored mode against $\mathcal{R}$)
    \State $D_{\mathsf{mint},i} \gets \mathcal{L}_i.T_0.D_\mathsf{mint}$;\quad ensure($D_{\mathsf{mint},i}.\mathsf{ty} = \mathcal{C}.\mathsf{ty}$)
    \State $(T_{\mathsf{b},i},u_{\mathsf{b},i},\pi_{\mathsf{b},i}) \gets$ final certified transfer of $\mathcal{L}_i$
    \State $(\predi_{\mathsf{b},i},h_{\mathsf{b},i}) \gets$ state spent by $T_{\mathsf{b},i}$ in $\mathcal{L}_i$
    \State ensure($T_{\mathsf{b},i}.\sthash = h_{\mathsf{b},i}$)
    \State $\mathsf{sid}_{\mathsf{b},i} \gets H(\predi_{\mathsf{b},i},\;T_{\mathsf{b},i}.\sthash)$;\quad $\txhash_{\mathsf{b},i} \gets H(T_{\mathsf{b},i}.D)$
    \State ensure($\Call{VerifyAnchoredLeaf}{\mathcal{T}.\alpha, \mathcal{R}, (\mathsf{sid}_{\mathsf{b},i},\txhash_{\mathsf{b},i}), \pi_{\mathsf{b},i}} = 1$)
    \State $\mathcal{R}_i \gets \Call{DecodeReason}{T_{\mathsf{b},i}.D.\auxd'}$ such that $T_{\mathsf{b},i}.D.\predi' = \predi_\mathsf{burn}(H(T_{\mathsf{b},i}.D.\auxd'))$
    \State \textbf{Destination:} ensure($\mathcal{R}_i.(\chi,a_\mathsf{V},a_\mathsf{A},\mathsf{ty},\mathsf{aid}) = (\chi,a_\mathsf{V},a_\mathsf{A},\mathsf{ty},\mathsf{aid})$)
    \State \textbf{Value:} $v_i \gets$ value of $D_{\mathsf{mint},i}.\auxd'$ for $\mathsf{aid}$ via $\Call{Assets}{D_{\mathsf{mint},i}.\mathsf{ty},D_{\mathsf{mint},i}.\auxd'}$
    \State ensure($\mathcal{R}_i.v = v_i$) \textbf{and} ensure($\mathcal{R}_i.v_\mathsf{f} \le v_i$)
    \State \textbf{Backing:} $(\,n_i,\mathcal{K}_i) \gets$ lock referenced by the genesis mint reason of $\mathcal{L}_i$
    \State $d_i \gets H(\mathsf{LOCK\_DOMAIN},\chi,a_\mathsf{V},n_i,\mathcal{K}_i)$
    \State ensure($\mathcal{K}_i.\mathsf{id} = \mathsf{id}(\mathcal{L}_i)$ \textbf{and} $\mathcal{K}_i.h_\mathsf{rcpt} = H(\mathsf{CBOR}(D_{\mathsf{mint},i}.\predi'))$)
    \State ensure($\mathcal{K}_i.\mathsf{aid} = \mathsf{aid}$ \textbf{and} $\mathcal{K}_i.v = v_i$)
    \State \textbf{Replay:} $\mathsf{btid}_{\mathsf{b},i} \gets H(\mathsf{BURN\_DOMAIN},\mathsf{sid}_{\mathsf{b},i},\txhash_{\mathsf{b},i})$;\quad $\eta_i \gets H(\mathsf{NUL\_DOMAIN},\mathsf{x}.h_\mathsf{cfg},\mathsf{btid}_{\mathsf{b},i})$
    \State ensure($\mathsf{VerifyNonMember}(A,\eta_i,w_i) = 1$);\quad $A \gets \mathsf{Insert}(A,\eta_i,w_i)$
    \State append $(\eta_i, \mathcal{R}_i.a_\mathsf{r}, v_i, \mathcal{R}_i.a_\mathsf{f}, \mathcal{R}_i.v_\mathsf{f}, \mathcal{R}_i.t_d)$ to $\mathsf{leaves}$
    \State append $(n_i,d_i)$ to $\mathsf{lref}$;\quad $V \gets V + v_i$
  \EndFor
  \State ensure($A = \mathsf{x}.A_\mathsf{new}$ \textbf{and} $V = \mathsf{x}.V$)
  \State ensure($\Call{Commit}{\mathsf{leaves}} = \mathsf{x}.\mathsf{rr}$)
  \State ensure($\Call{Commit}{\Call{DedupSort}{\mathsf{lref}}} = \mathsf{x}.\mathsf{lr}$)
  \State \Return $1$
\EndFunction
\end{algorithmic}
\end{algorithm}

The single source lock per burn shown above is the supported case. A vector of source locks per burn is reserved for a future merge/combine primitive; until its verifier is defined, a relation MUST reject any burn that references more than one lock.

\section{Settlement}\label{app:bridging-settlement}

The vault keeps: the verification key $\mathsf{vk}$, the configuration hash $h_\mathsf{cfg}$, the custody asset $a_\mathsf{A}$, an allow-list of trust-base hashes (rotatable under timelock for validator-set epochs), the bridge-in map $\mathsf{lockDigest}[\cdot]$, and the accumulator root $A$ (initialized to $A_\emptyset$). Algorithm~\ref{alg:bridge-settle} settles a batch.

\begin{algorithm}[!htbp]
\caption{Settlement $\textsc{FulfilBatch}$}\label{alg:bridge-settle}
\begin{algorithmic}[0]
\Function{FulfilBatch}{$\mathsf{x},\Pi,\langle\ell_i\rangle,\langle(n_j,d_j)\rangle$}
  \State ensure($\mathsf{Verify}(\mathsf{vk},\mathsf{x},\Pi) = 1$)
  \State ensure($\mathsf{x}.h_\mathsf{cfg} = h_\mathsf{cfg}$ \textbf{and} $\mathsf{x}.h_\mathcal{T} \in \mathsf{trustBaseAllowed}$)
  \State ensure($\mathsf{x}.A_\mathsf{old} = A$) \Comment{the entire replay guard}
  \State ensure($\len{\langle\ell_i\rangle} = \mathsf{x}.B$ \textbf{and} $\Call{Commit}{\langle\ell_i\rangle} = \mathsf{x}.\mathsf{rr}$)
  \State ensure($\langle(n_j,d_j)\rangle$ sorted-unique by $n_j$ \textbf{and} $\Call{Commit}{\langle(n_j,d_j)\rangle} = \mathsf{x}.\mathsf{lr}$)
  \ForAll{$(n_j,d_j)$}
    \State ensure($\mathsf{lockDigest}[n_j] = d_j$)
  \EndFor
  \State ensure($\sum_i \ell_i.v = \mathsf{x}.V$ \textbf{and} $\forall i:\ \ell_i.v_\mathsf{f} \le \ell_i.v$)
  \State $A \gets \mathsf{x}.A_\mathsf{new}$ \Comment{one root update for the whole batch}
  \For{$i \gets 1 \ldots \mathsf{x}.B$}
    \State $f \gets (\mathsf{now} \le \ell_i.t_d \textbf{ and } \ell_i.a_\mathsf{f} \ne \bot)\ ?\ \ell_i.v_\mathsf{f}\ :\ 0$ \Comment{deadline and recipient gate the fee only}
    \State transfer $\ell_i.v - f$ of $a_\mathsf{A}$ to $\ell_i.a_\mathsf{r}$;\quad \textbf{if} $f \ne 0$ transfer $f$ to $\ell_i.a_\mathsf{f}$
    \State emit $\mathsf{Released}(\ell_i.\eta, \ell_i.a_\mathsf{r}, \ell_i.v, f)$
  \EndFor
  \State emit $\mathsf{BatchFulfilled}(\mathsf{x}.A_\mathsf{old}, \mathsf{x}.A_\mathsf{new}, \mathsf{x}.\mathsf{rr}, \mathsf{x}.B)$
\EndFunction
\end{algorithmic}
\end{algorithm}

Properties:
\begin{itemize}
  \item \textbf{No external query, no light client.} The proof self-certifies its anchor under the allow-listed $\mathcal{T}$, and backing is confirmed against locally recorded digests; the vault tracks no Unicity roots.
  \item \textbf{Fee optionality never strands funds.} $t_d$ and $a_\mathsf{f}$ gate only the fee. After the deadline, or when $a_\mathsf{f}=\bot$, the recipient receives the full amount; anyone (including the recipient) can submit the proof, so a slow or absent relayer costs only the fee.
  \item \textbf{Batch atomicity.} All $B$ transfers settle together; a single recipient that rejects receipt reverts the batch. This is a liveness concern, not a safety one: on revert $A$ is unchanged and no nullifier is consumed, so the offending leaf is dropped and the batch rebuilt. A deployment expecting blockable recipients\footnote{e.g., compliance reasons} may settle by crediting a withdrawable balance instead of directly transferring; doing so also isolates one misbehaving or reverting recipient from every other release leaf in the same batch, which matters because a proof, once produced, commits an entire batch as a unit.
  \item \textbf{Reconstructability.} $\mathsf{Released}$ emits each nullifier and $\mathsf{BatchFulfilled}$ each root transition, so the off-chain accumulator is rebuildable from public data alone and verified against the on-chain $A$ before use.
  \item \textbf{Transfer gas/energy.} On EVM-class chains, a low-level call the vault issues to move the asset to $\ell_i.a_\mathsf{r}$ does not necessarily receive the caller's full remaining gas or energy by default, particularly immediately following an expensive operation such as proof verification within the same transaction; some execution environments cap or otherwise constrain forwarding to nested calls. A vault implementation should budget an explicit, sufficient tx fee for each transfer rather than relying on default forwarding.
\end{itemize}

\section{Security Properties}\label{app:bridging-security}

\begin{description}
\item[Solvency.] By bridge-in (Sec.~\ref{app:bridging-backing}), every bridged token is backed by an equal lock; by split verification (Sec.~\ref{sec:split-verification}), the value of any split descendant is at most the parent's. The backing check ties each burn to a recorded lock digest, and value conservation ties the burned amount to that lock's amount, so $\sum(\text{released}) \le \sum(\text{locked})$ and custody cannot be over-drawn.
\item[Replay-freeness.] Nullifiers are globally unique (Sec.~\ref{app:bridging-nullifier}). Each release proves non-membership and inserts, and the vault requires $\mathsf{x}.A_\mathsf{old} = A$. A resubmitted proof fails (the stored root advanced); a fresh proof cannot show non-membership of an inserted nullifier; duplicate nullifiers within a batch fail the second non-membership check.
\item[Authenticity and finality.] $\mathfrak{R}$ embeds $\Call{VerifyToken}{}$ against the pinned $\mathcal{T}$; by knowledge soundness, an accepting proof implies a genuine, BFT-final burn whose reason is exactly $\mathcal{R}$. The recipient and amount cannot be altered, as they are committed inside the certified burn.
\item[Liveness.] The prover holds no authority and is replaceable: from the on-chain root and the public $\mathsf{Released}$ stream, any party rebuilds the accumulator, refetches anchored inclusion proofs from the Unicity Service, and reproduces a proof. The residual requirement is availability of the burned-token data, which the holder must retain until settlement.
\end{description}

\section{Batched Proving}\label{app:bridging-batching}

The public statement and the settlement procedure are identical for every value of $B$, so the vault and the wallet are stable across the following stages; only the prover changes.

\begin{description}
  \item[Single ($B=1$).] One token, one nullifier insertion, one root transition. Establishes the interface.
  \item[In-circuit batch ($B>1$).] One proof for $B$ burns sharing one anchor seal $C^{\mathsf{r}*}$ and one accumulator transition. This amortizes, over the whole batch, the single on-chain $\mathsf{Verify}$, the single seal verification, and the single root update. The remaining on-chain settlement work is linear in the number of release leaves, with asset transfers the irreducible dominant term; lock checks are linear in the number of distinct referenced locks.
  \item[Recursion.] Prove each token's lineage independently (parallel across machines), each emitting its nullifier and validated facts; an aggregation relation verifies the child proofs, performs the batch accumulator insertion, binds all returns to one $h_\mathsf{cfg}$ and trust-base policy, and emits the same statement $\mathsf{x}$. This decouples per-token proving latency from batch size.
\end{description}

\section{Relation and Configuration Upgrades}\label{app:bridging-upgrades}

Both $\mathsf{vk}$ and $h_\mathsf{cfg}$ are fixed at vault deployment (Sec.~\ref{app:bridging-config}, \ref{app:bridging-primitives}) and neither can be changed in place. $h_\mathsf{cfg}$ is in fact self-referential through $a_\mathsf{V}$: a new vault address necessarily produces a new $h_\mathsf{cfg}$, so relocating custody to a new vault instance and changing the proving relation are, mechanically, the same kind of event even when only one of them was the actual motivation.

A relation upgrade, e.g., correcting a defect in $\mathfrak{R}$ (Alg.~\ref{alg:bridge-relation}) or extending it, proceeds as:
\begin{enumerate}
  \item Deploy a new vault with the new $\mathsf{vk}$ and the resulting $h_\mathsf{cfg}$, initialized with an empty accumulator ($A = A_\emptyset$), an empty lock-digest map, and no allow-listed trust base; each MUST be populated before the vault is usable.
  \item Publish the new configuration to wallets and provers, and stop bridge-in against the superseded vault.
  \item Return outstanding value locked in the superseded vault under the \emph{old} relation. The new vault has no record of old locks and cannot settle proofs referencing them, so this step cannot be deferred to the new deployment.
\end{enumerate}
Value not returned before the superseded vault is abandoned is unrecoverable through the protocol: $\mathsf{lockDigest}$ (Sec.~\ref{app:bridging-lock-op}) is local, per-instance state.

This rigidity is the direct consequence of $\mathsf{vk}$ being immutable, which is what makes the vault trustworthy in the first place.

\section{Shared Parameters}\label{app:bridging-params}

The table below fixes the bridge's free parameters with initial defaults. The \textbf{Status} of each entry reads: \emph{Fixed} --- determined by this specification; \emph{Default} --- a RECOMMENDED value a deployment SHOULD adopt unless it has a specific reason otherwise; \emph{Rule} --- a mandatory derivation constraint; \emph{Chain} / \emph{Policy} --- deployment- or chain-specific, to be chosen and published before deployment of the source chain contracts. Every value, once chosen, MUST be frozen and published so that wallets, provers, and the vault agree.

\begin{longtable}{p{3.0cm}p{4.2cm}p{2.0cm}p{3.4cm}}
\textbf{Parameter} & \textbf{Default / value} & \textbf{Status} & \textbf{Notes} \\
\hline
\endhead
\multicolumn{4}{l}{\textit{Hashing and identifiers}}\\
\hline
$H$ & $\mathsf{SHA\text{-}256}$ over deterministic CBOR & Fixed & all bridge hashes and digests \\
$\mathsf{ver}$ & $1$ & Fixed & version of backing and return reasons \\
\hline
\multicolumn{4}{l}{\textit{Domain separators and CBOR tags}}\\
\hline
$\mathsf{CFG\_DOMAIN}$    & \texttt{"UNICITY\_BR\_CFG"}  & Default & separator in $h_\mathsf{cfg}$ \\
$\mathsf{LOCK\_DOMAIN}$   & \texttt{"UNICITY\_BR\_LOCK"} & Default & separator in lock digest $d$ \\
$\mathsf{NUL\_DOMAIN}$    & \texttt{"UNICITY\_BR\_NUL"}  & Default & separator in nullifier $\eta$ \\
% $\mathsf{BACKING\_TAG}$   & CBOR tag \texttt{390xx}        & Default & genesis backing reason \\
$\mathsf{REASON\_TAG}$    & CBOR tag \texttt{39048}        & Default & return reason $\mathcal{R}$ \\
\hline
\multicolumn{4}{l}{\textit{Encodings}}\\
\hline
Amounts $v,v_\mathsf{f}$ & unsigned, $<2^{256}$, minimal big-endian, $\le 32$ bytes & Fixed & as split amounts (Sec.~\ref{app:rsmst}) \\
Addresses $a_\mathsf{V},a_\mathsf{A},a_\mathsf{r},a_\mathsf{f}$ & $20$ bytes & Chain & EVM-class default; set by $\chi$ \\
$\mathcal{C},\mathcal{K},\mathcal{R}$, backing reason, lock event, release leaf, lock ref & fixed-width CBOR arrays, stated field order & Fixed & deterministic CBOR, consumed completely \\
\hline
\multicolumn{4}{l}{\textit{Derivations (mandatory)}}\\
\hline
$h_\mathsf{cfg}$ and $a_\mathsf{A}$ & jointly from one config & Rule & Sec.~\ref{app:bridging-config} \\
Nullifier $\eta$ & per Sec.~\ref{app:bridging-nullifier} & Rule & binds $h_\mathsf{cfg}$ and the certified burn leaf \\
Lock digest $d$ & per Sec.~\ref{app:bridging-lock-op} & Rule & vault-stored at bridge-in \\
Batch witness order & per Sec.~\ref{app:bridging-primitives} & Rule & witness $k$ vs.\ root after inserts $0..k{-}1$ \\
\hline
\multicolumn{4}{l}{\textit{Nullifier accumulator}}\\
\hline
Tree type & radix sparse Merkle tree (Sec.~\ref{app:rsmt}) & Default & native non-inclusion; key $=\eta$ \\
Key space & $256$-bit, sparse & Fixed & $\eta \in \bytes{32}$ \\
Leaf value & fixed presence constant \texttt{0x01} & Default & only membership is needed \\
Accumulator hash & $\mathsf{SHA\text{-}256}$ & Default & matches the Service SMT \\
Empty root $A_\emptyset$ & empty-RSMT root ($\zerohash$) & Fixed & vault initialization \\
\hline
\multicolumn{4}{l}{\textit{Calldata commitment}}\\
\hline
$\Call{Commit}{}$ & deployment-fixed binding vector commitment & Rule & pinned by $\mathsf{vk}$ and the vault; EVM-class deployments use keccak-256 over fixed-width ABI encodings \\
\hline
\multicolumn{4}{l}{\textit{External finality}}\\
\hline
$\Call{LockFinal}{}$ realization & live query with $K$ confirmations & Chain & or carried inclusion proof vs.\ pinned $\chi$ anchor (Sec.~\ref{app:bridging-finality}) \\
$K$ (confirmation depth) & $20$ (Tron-class); ``finalized'' (Ethereum L1) & Chain & set for $\chi$'s reorg resistance \\
\hline
\multicolumn{4}{l}{\textit{Batching and settlement}}\\
\hline
$B_{\max}$ (max batch size) & $64$ & Tunable & bounded by external block gas/energy \\
Deadline $t_d$ & gates the fee only; principal always released & Fixed & Sec.~\ref{app:bridging-settlement} \\
Fee constraint & $v_\mathsf{f} \le v$ & Fixed & enforced in circuit and vault \\
\hline
\multicolumn{4}{l}{\textit{Governance and proof system}}\\
\hline
Trust-base allow-list & timelocked add/rotate & Policy & for validator-set epochs \\
Governance timelock $\Delta_\mathsf{gov}$ & $\ge 7$ days (suggested) & Policy & deployment-set \\
Proof system & Groth16-wrapped recursive STARK & Default & Wrapping applies to this boundary only (Sec.~\ref{app:bridging-primitives}); $\mathsf{vk}$ pinned per circuit version, and changing it is a vault redeployment (Appendix~\ref{app:bridging-upgrades}) \\
\end{longtable}

A deployment must select the chain-specific entries ($\chi$, the address width, $K$ or the external anchor), pinned $\mathsf{vk}$ for the chosen proof system, and publish the resulting $h_\mathsf{cfg}$ and the allowed trust-base hashes.
